 \section{Pyraf}

Currently the reduction package is running in a docker file.

\begin{tcolorbox}[colframe=red!75!white,colupper=black,title=Pyraf alias entry for bash]
\begin{verbatim}
function pyraf { export LOCALDATA=$(pwd);
                 export SPECDATA=/mnt/MySharedDir;
                 cd /home/git/external/PyrafDocker/Compose/pyraf;
                 sudo xhost +local:docker;
                 docker compose run --remove-orphans pyraf /bin/bash;
               }
\end{verbatim}
\end{tcolorbox}

Where: LOCALDATA is a directory on the host, SPECDATA may be /mnt/MySharedDir.
This gives you the ability from within the docker, to copy files from
a remote machine to the host machine for faster consideration.\index{PyRAF!shared directory}
\index{shared directory}.

This CDs to a local git repository where the \dhl{docker-compose.yml} file lives,
and starts the docker container running. Note: The \dhl{sudo xhost +local:docker}
statement allows the container to use the local X server.

\ltodo{Win11}{Update for Win11 docker use.}

The new PyRAF community edition has made significant changes to the
overall architecture.

The \dhl{docker-compose.yml} file ties several directories to gether.
You may want to alter those locations for your situation.





%% 229, 224, 213
%%  (mapcar '(lambda (a) (/ a 255.0)) (list 229 224 213))
%% (0.8980392156862745 0.8784313725490196 0.8352941176470589)


